and related theory suggests that phenology may determine community assembly \citep{wolkovich2021phenological,levine2022competition}.'



The demonstrated importance of seasonal priority effects in mediating interspecific competition has led to growing interest in the role of seasonal priority effects in coexistence \citep{Diez_2024, Blackford_2020}. Several existing community assembly models that integrate the timing of species arrival or activity suggest timing is fundamental to coexistence \citep{volkerisaj}, with multiple studies proposing a fundamental trade-off between early phenology to gain unrestricted access to  resources versus later phenology, which requires competing well under low resource conditions \citep{wolkovich2021phenological,levine2022competition}. Growing interest in understanding seasonal priority effects also highlights one of the major challenges to the field: phenology is not a static trait. Phenology varies highly across years because of the contribution of the environment in determining the timing of event for most, if not all, species. 


The importance of seasonal priority effects in mediating interspecific competition suggests how phenology may be fundamental to coexistence \citep{Diez_2024, Blackford_2020}, and thus also to forecasting community change with warming. Functionally, phenology is the realized product of an interaction between the environment (cues) and species unique physiology, genetic and developmental attributes \citep{DONOHUE2005}. 



These differences are likely amplified by climate change

observed in different environmental conditions (e.g., in observational studies across years). Realized phenology differences between species can substantially depending on the environmental conditions.

Leveraging community assembly models to address how shifting phenology with climate change may affect coexistence requires models that encapsulate the strong inter-dependence between environmental variability to competition strength via phenology.

% If timing determines community assembly, then shifts in the environment could lead to local species extinctions


Comparing long-term co-occurrence (which we consider coexistence here) of species pairs under two different climate regimes we found that each required totally different trait variation such that no pair co-existing under one climate could also co-exist under the other (Figure \ref{Fig:differences}a, see circled points for example, results from 1000 runs of 300-year duration under two different climate normals; see Extended Method for details). 


This is because the relative phenological differences between species that structures the tradeoff are dependent on the interaction between each species' phenological sensitivity and the environment itself, and therefore, in a novel environment the relationship between differences in species phenological sensitivities and realized phenological differences shift (Figure \ref{Fig:differences}b). 


Our model was similar to others in varying species intrinsic competitive abilities at the same time we varied phenology, but differed in modeling phenological sensitivities, making phenology each year a direct response to the intra-annual environmental variation. % Megan does not like 'phenology each year a direct response to the intra-annual environmental variation' because in our model we draw chilling once per year, it's not intra-annually modeled. Intra-annually we model what happens AFTER the SOS (after chilling is over). 


Our modeling extension revealed an additional, important point for improving the predictive power of coexistence model with climate change. While the tradeoff between realized phenological differences and competitive differences is stable across a shifting environment (i.e., the same magnitude if phenological differences are required to offset the same level of differences in competitive ability), it is the tradeoff between phenological sensitivity differences and competition differences that change in a novel environment (Figure \ref{Fig:differences}a,c). Thus, measuring species phenologies may far less useful for understanding future community assembly than measuring their phenological sensitivities.

Our model shows that the same trade-off between realized phenological differences and competitive differences across environments 

In a rapidly changing world, these kind of cross-disciplinary collaborations are critical to translate ecological theory into practice and advance a more predictive ecology than is better equipped to inform ecosystem conservation and management.

Paths forward...

Estimates of species-level phenological sensitivity, quantified at multiple scales (i.e., population, individual) would allow researchers to parameterize community assembly models for species of interest. This could form the start of practical application of these models, building the groundwork for modelers and experimentalists to work together to build predictive tools that aid ecological practitioners. But this will require expanding the study systems and flexibility of current models by both research groups. 

Models of ecological competition that include phenology are often developed to encapsulate competition between two annual plants \citep{godoy2014,wolkovich2021phenological,levine2022competition}....
To help, modelers can construct more biologically realistic community assembly models that include multiple growth forms and life stages. At the same time, experimentalists could design studies that explore how environmental variability influences phenological differences across life-stages and mediates both intra- an inter-specific competition within and across these stages.


% emw and MJD 15Feb2026 -- suggest to cut
% Another dimension of risk and phenology could help generate more predictive assembly models is to understand how another aspect of phenology that is sensitive to environmental cues---the duration of a phenological phases---mitigates or exacerbates environmental risk. Environmental cues not only advance or delay the start of a phenological transition, but may also influence the rate of the transition. More protracted phenological phases (i.e., if the germination of a full seed cohort takes longer) associated with unfavorable conditions may serve as a form of within-season bet-hedging in which an untimely mortality event would have a lower fitness cost to a cohort of seed that started germinating early but slowly (at the cohort level) than one that started early and finished fast.%emw9Aug2025:  Could cite some eco-evo climate change papers here?